\documentclass[../../main2.tex]{subfiles}

\begin{document}

	% 全书位置：导论(2008) → 【动机】 → 目标 → 预测与因果 → ……
	% v2 定位（notes/02-architecture-v2.md §8）：动机不再是「解释的起点」，
	% 而是「预测的必要输入」——2008 年的模型输就输在缺这一味。

	\textbf{\large 全书位置}：导论 → \textbf{动机} → 目标。

	\textit{把「他这么做，是因为……」这句话拆到不能再拆，找出预测社会的模型里最不可省的那一味输入。}

	\vspace{1em}

	\chapter{动机：解释的最小起点}
	\label{ch:motivation}

	\begin{motivation}
		导论留下的诊断：华尔街的模型输掉 2008，不是因为数据少、算力小，而是因为模型里缺一味输入——局中人的动机。缺了它，看得见「价格一起动」，看不见「谁在推」。\\
		\textbf{核心动机}：把「他这么做，是因为……」这个日常动作拆到最底层，论证为什么\textbf{任何想预测社会的模型，都必须把「动机」当必要输入}——它不可观测，但不可省略。
	\end{motivation}

	\begin{logicmap}
	\begin{tikzpicture}[node distance=0.9cm, every node/.style={draw, rectangle, rounded corners, align=center}]
	\node (q) {「他这么做，是因为……」};
	\node (unit) [below=of q] {最小单元：主体+动机+环境};
	\node (debris) [below left=0.9cm and 1.4cm of unit] {残骸 vs 蜂巢};
	\node (order) [below=1.6cm of unit] {动机先于结构};
	\node (res) [below=of order] {动机 $\neq$ 结果};
	\node (div) [below right=0.9cm and 1.4cm of res] {同一动机·两种社会};
	\node (next) [below=1.6cm of res] {目标？（抛给第2章）};
	\draw[-{Stealth}] (q) -- (unit);
	\draw[-{Stealth}] (unit) -- (debris);
	\draw[-{Stealth}] (unit) -- (order);
	\draw[-{Stealth}] (debris) -- (order);
	\draw[-{Stealth}] (order) -- (res);
	\draw[-{Stealth}] (res) -- (div);
	\draw[-{Stealth}] (res) -- (next);
	\draw[-{Stealth}] (div) -- (next);
	\end{tikzpicture}
	\end{logicmap}

	% ==================================================================
	\section{一个最小的解释单元}
	\label{sec:mot-minimal}

	\begin{readingnote}
		你有没有想过，为什么小孩问「为什么」总能问到大人哑口无言？\\
		大人不是无知——是任何解释都会走到尽头。尽头处站着的那个人，是谁？
	\end{readingnote}

	先看一个再日常不过的场景。小孩打碎了杯子，大人问：为什么碎了？——因为手滑。为什么会手滑？——因为杯子烫。为什么杯子烫？——因为刚倒了热水。为什么倒热水？——因为想喝茶。为什么想喝茶？

	问到第四个「为什么」的时候，大人开始不耐烦了。注意这个不耐烦的时刻：它不是耐心的极限，是\textbf{解释的极限}。再往下答，大人给出的已经不是「喝茶的原因」，而是神经递质、血糖曲线、演化史——那些是别的学科的地盘。解释链必须停在某处；停下的地方，就是这条解释链的\textbf{起点}。

	本书的主张是：对社会现象，最经济的停点是\textbf{动机（motivation）}。想喝茶——一个动机。它不可再拆，不是说它没有生理基础，而是说再拆下去，你就从「解释社会」滑进了「解释神经系统」，问题换了，答案自然对不上。

	现在换上预测者的眼镜，把「他这么做，是因为……」这句话正式拆开。它一共用到三件东西：
	\begin{itemize}
		\item \textbf{主体}：「他」——一个会行动的单位（人、组织、乃至一家投行）；
		\item \textbf{动机}：「因为」后面接的那个东西——主体想要什么；
		\item \textbf{环境}：这句话没说出口、但永远默认存在的背景——杯子是烫的、水是有的、茶是能买到的。
	\end{itemize}

	任何解释都需要这三件套；而任何\textbf{预测}，需要的是同三件的未来版本：他还会是那个他吗？他还想要那个东西吗？环境还惯着他吗？2008 年的模型三缺二：它精细地刻画了环境（利率、房价、违约率的历史），把主体压成一个统计平均数，把动机整个留白。三缺二的模型可以精确描述过去，不可能预测未来。

	\begin{questionbox}
		\textit{现在你可能想反驳：「动机是最小单元？那动机本身又是从哪来的？你只是把问号往下挪了一层。」}\\
		\textbf{诚实回答：本书不假装能回答这个问题。}动机的来源——演化、生理、童年、历史——在本书的边界之外。本书从动机开始，就像几何学从「点」开始：几何学不解释点是什么做的，只规定点服从什么规则。我们也一样：不解释动机从哪来，只研究动机服从什么规律。至于「环境塑造下一轮的动机」——那不是对起点的追问，而是循环的闭合，第\ref{ch:closure}章会回来。
	\end{questionbox}

	\begin{reader_anticipation}
		\item 动机看不见摸不着，怎么研究？→ 不能直接观测，只能从行为反推——第2章「显示性偏好」就是干这个的（名字到那章再正式给）。
		\item 为什么把「环境」也算进最小单元？它看起来只是背景板。→ 因为第1.3节将看到：同样的动机，环境不同，结果相反。环境不是背景板，是乘数。
	\end{reader_anticipation}

	\paragraph*{逻辑衔接}
	我们有了最小单元：主体、动机、环境。但马上冒出一个次序问题——这三件里，凭什么把「动机」放在最前面？结构、环境难道不比动机更「硬」、更不需要解释吗？下一节正面回答：为什么动机先于结构。

	% ==================================================================
	\section{动机为什么先于结构}
	\label{sec:mot-structure}

	\begin{readingnote}
		你有没有想过，为什么我们看到蜂巢会觉得「精巧」，看到地震后的碎砖堆只觉得「乱」？\\
		两者都是物质的排列——差别到底在哪？这两个画面里，是不是藏着\textbf{「为了什么」的在场与缺席}？
	\end{readingnote}

	直觉先给答案：精巧与乱的区别，不在几何。完全可以想象一堆碎砖，其轮廓恰好酷似一座蜂巢——形状相同，一个让你觉得「有目的」，一个让你觉得「没目的」。差别在于：蜂巢的形状是\textbf{被一个目标筛出来的}，碎砖的形状只是力学结果。

	地震后的墙壁残骸就是一个标准的\textbf{无目标的结构}：它有形状、有内部排列、甚至有「结构」这个词汇的一切表面特征。但解释它不需要「想要」——问「残骸为什么是这个形状」，答案全在受力与材料，一句「为了什么」都用不上。

	现在用本书的标准句式，把「动机」从这个对照里分离出来：

	\begin{quote}
		如果一个结构不需要「为了什么」就能被解释，那它就是废墟——性质是纯力学的，预测它只需要物理。但现在加入一个约束：\textbf{这个结构必须持续存在}——持续存在需要主动对抗消耗：修复、复制、防御。这一步改变了所有性质，这就是「动机」从「结构」中分离出来的时刻。
	\end{quote}

	为什么「持续存在」需要一个筛选者？这里把前身第1章的论证压缩成一段。\textbf{存在是小概率事件}：精确的环境参数、稳定的能量输入、多个随机事件的巧合叠加——每个条件都是小概率，凑在一起近乎奇迹。这就像中彩票：选对一个数字很难，连续选对七次，难到几乎不可能。

	既然存在如此偶然，任何「还在」的结构就都自带一个倾向：\textbf{继续存在}。这个倾向不是设计出来的，是被筛选出来的——存在过但不维持的东西，我们已经看不见了。你之所以能读到这一页，是因为你的每一个祖先都成功地没有死绝。我们觉得「存在」遍地都是，只是因为只有存在者才能提问——这是观察者选择效应（observer selection bias）。

	于是自然选择给出了「动机」的第一个形式化：它不是选择最优，而是\textbf{淘汰不可行}。过滤器只问一个问题——「还能维持吗？」——凡不能维持的，从解释对象里被删除。这个「维持」的倾向，就是动机的最原始形式：不需要意识，不需要计划，只需要\textbf{还站着}。

	\lowfreq{注意：把「维持存在」称为动机，是一个建模约定，不是定理。}它的好处是让「石头没有动机、细菌有动机、公司有动机」这件事可以被同一套语言处理；它的代价是把「动机」从心理学词汇降格为系统论词汇。本书接受这个代价，因为我们要的是跨层级的解释力。

	排序到此清楚了：\textbf{先有「维持存在」这个筛选标准，后有留存下来的结构}。结构不是初始事实，是筛选后的稳态。所以解释一个结构，必须先给出筛选标准（它为了什么被留下），再解释为什么恰好是这一个被留下（环境与历史）。预测也一样：想知道一个组织会不会僵化、一家银行会不会冒进，不能只看它今天的结构——要看它还被什么目标筛着。2008 年之前，华尔街的结构（评级机构、奖金制度）运转良好，恰恰因为它们还筛在「短期利润」这个目标上；目标没变，结构不可能自己变。

	\begin{questionbox}
		\textit{现在你可能想反驳：「先讲结构不行吗？结构主义、地理决定论，都是先讲结构的成功范式。」}\\
		\textbf{可以，但只做半件事。}先讲结构能回答「什么结构稳定」，答不了「为什么会有结构」——地形不想要任何东西，想要东西的是地形上的人。地理决定论解释约束，但不解释这些约束为什么被这样利用。所以在本书的新结构里，地理被降级安置：它是「看得见的约束」——第2章的一节工具，不再是原因的终点。
	\end{questionbox}

	\begin{reader_anticipation}
		\item 地理决定论在新书里去哪了？→ 第2章 §「看得见的约束」：降级为工具，但没有被丢弃。
		\item 「淘汰不可行」的过滤器，有正式的形态吗？→ 第\ref{ch:channels}章「通道生成器」：有压扩增与无压扩增，就是过滤器的两种数学形态。
	\end{reader_anticipation}

	\paragraph*{逻辑衔接}
	动机拿到了起点的位置，结构退居筛选后的结果。但读者应该已经闻到一丝危险：如果动机是起点，动机能保证结果吗？下一节给出坏消息——不能。而这个坏消息，正是这本书剩下的全部篇幅存在的原因。

	% ==================================================================
	\section{动机不决定结果}
	\label{sec:mot-result}

	\begin{readingnote}
		你有没有想过，为什么两个都想活下去、都想让家人吃饱的人，一个组织起来修了水渠，一个造了商船——最后活成了完全不同的社会？\\
		同一个动机，怎么会长出两个世界？这背后是不是有\textbf{动机和结果之间的某种断裂}？
	\end{readingnote}

	把镜头拉到两千多年前。黄河流域的农夫和爱琴海的渔民，动机清单几乎一字不差：安全、吃饱、繁衍、免于暴力。谁也不比谁更「爱好集权」或「热爱自由」。

	但环境把同样的动机揉出了相反的形状。黄河年年泛滥，治水需要跨村落的协作，协作需要集中调度——在这个环境里，集权的成本被治水的收益覆盖了，于是\textbf{大一统}被筛出来。爱琴海多山破碎、岛屿密布，农耕分散而贸易有利可图，贸易需要的是契约和港口而不是调度——在这个环境里，分权的成本被贸易的收益覆盖，于是\textbf{城邦}被筛出来。

	动机相同，环境不同，结果相反。这就是本节要钉死的一句话：\textbf{动机给方向，环境给形状；结果是动机和环境相乘，不是动机的延长线}。

	\textit{生活类比}：同样的种子，撒在湿地长成芦苇，撒在岩缝长成针叶——你不能指着芦苇说「种子想要潮湿」。（诚实标注：生态学里这套机制的真实名字是\textbf{生态位分化（niche differentiation）}——同一物种在不同环境下长出不同形态。这里是它的通俗版，当类比读，别当论证用。）

	现在，请注意刚才那两段话里偷偷多出来的一个东西。农夫并不天生想修渠——他想要的是\textbf{不被淹}。修渠只是手段。于是「想要」其实分了三层：
	\begin{itemize}
		\item \textbf{动机}：想要什么（不被淹、活下去）；
		\item \textbf{目标}：把想要的东西定成一个可比较的对象（「五年内洪水不进村」）；
		\item \textbf{结果}：世界实际给了什么（一座渠、一套官僚制、一个王朝）。
	\end{itemize}

	从动机到结果，隔着两层转换：\textbf{目标化}（把愿望变成对象）和\textbf{环境化}（把对象放进世界）。动机不决定结果，因为它管不了这两层。

	这里还藏着一个对预测者更坏的消息：从结果往回推，是\textbf{推不回来的}。同一个结果可以由不同的动机配不同的环境造成——大一统可以来自治水，也可以来自御敌，还可以来自盐铁专营；只看「大一统」这个结果，你还原不出是哪一对组合。回溯式的解释永远丢信息。这不是态度问题，是方向问题：想从「结果」走到「原因」，路是断的；想从「原因」走到「结果」，路才是通的。整本书走的是后一条路。

	\begin{questionbox}
		\textit{现在你可能想反驳：「既然动机不决定结果，那要动机干什么？直接从环境推结果不就行了？地理决定论干脆复活算了。」}\\
		\textbf{不行，反例同款：同一环境下也能有不同结果。}同一条河，两岸制度可以不同；同一场危机，邻国应对可以相反。环境单独也不充分。两个都不充分，就说明预测必须同时带三样东西：\textbf{看得见的约束}（环境，第2章）、\textbf{藏在行为里的目标}（第2章）、以及\textbf{两者互相推挤的方式}——推挤的产物，就是第\ref{ch:channels}章以后要拆的「齿轮」。本书不搞任何单因决定论——包括以「动机」为单因的版本。
	\end{questionbox}

	\begin{reader_anticipation}
		\item 动机到结果之间那道「落差」，用什么衡量？→ 第2章：目标必须先成为可比较的量。
		\item 「同一环境下不同结果」，「相关」和「原因」怎么分？→ 第\ref{ch:prediction}章：数据的第一个陷阱。
		\item 多个因素共同导致一个结果，「各占多少」怎么算？→ 第\ref{ch:contribution}章「无你检验」——本书的定量核心。
	\end{reader_anticipation}

	\paragraph*{逻辑衔接}
	我们手里现在有一件麻烦事：动机是起点，但起点不保证终点；动机和结果之间，隔着「目标」和「环境」两道转换，而且从结果回推的路是断的。环境这一维，下一章以「看得见的约束」登场；眼下最紧迫的是中间那道——「目标」。农夫想要的东西，到底是个什么东西？它有什么形状？下一章拆这个。

	% ==================================================================
	\section{本章小结：目标必须先被规定}
	\label{sec:mot-summary}

	\begin{readingnote}
		你有没有想过，为什么「我知道他想要什么」这句话说出口之后，我们还是预测不了他会做什么？\\
		因为我们其实还不知道，他想要的东西\textbf{长什么样}——是一个名词、一张清单，还是一个数？
	\end{readingnote}

	收束本章的三步：
	\begin{enumerate}
		\item \textbf{最小单元}（§\ref{sec:mot-minimal}）：任何解释、任何预测，都建立在主体+动机+环境的三件套上；解释链必须停，对社会现象最经济的停点是动机。缺了动机的模型，就是 2008 年那个三缺二的模型；
		\item \textbf{次序}（§\ref{sec:mot-structure}）：动机先于结构，因为结构是「维持存在」这个筛选标准过滤后的稳态——先有标准，后有幸存者；
		\item \textbf{限度}（§\ref{sec:mot-result}）：动机不决定结果；从动机到结果隔着目标化与环境化两层转换，回溯必丢信息。
	\end{enumerate}

	下一步的路因此是明确的：要谈「动机与结果之间的落差」，必须先把「目标」规定成能比较、能排序的东西——否则「落差」只是修辞。说一个人「想要 X」，X 到底是什么形状？这个问题就是第2章的开门钥匙。

	我们能说什么：预测的最小原料是三件套，且原料内有序；动机是筛选后的倾向，结构是筛选后的结果。我们不能说什么：动机的最终来源；从动机到结果的确定性预测——那需要等目标（第2章）、因果（第\ref{ch:prediction}章）和无你检验（第\ref{ch:contribution}章）全部就位。

	\paragraph*{逻辑衔接}
	本章把「动机」安放在了预测模型的必要输入位上，并诚实地交出了它的局限。但「目标」此刻还是一个日常词汇——日常词汇没法计算。第2章「目标：把意向写成可比较的量」将做这件事：让「想要」变成能排序、能打分、能从行为里反推出来的东西。

\end{document}
